\subsection{Đề 3}
\Opensolutionfile{ans}[ans/ans-2-OHK2-De3]
\caulc
%nlc 1
\begin{ex}%[2D4H1-2] 
	Biết rằng $\displaystyle\int\limits_{0}^{2}f(x)\mathrm{\,d}x=-4$. Giá trị của $\displaystyle\int\limits_{0}^{2}[3x-2f(x)]\mathrm{\,d}x$ bằng
	\choice
	{$-2$}
	{$12$}
	{\True $14$}
	{$22$}
	\loigiai{
		Ta có $\displaystyle\int\limits_{0}^{2}[3x-2f(x)]\mathrm{\,d}x=\displaystyle\int\limits_{0}^{2}3x\mathrm{\,d}x-2\displaystyle\int\limits_{0}^{2}f(x)\mathrm{\,d}x=\left(\dfrac{3x^2}{2}\right)\bigg|_{0}^{2}-2(-4)=6+8=14$
	}
\end{ex}
%nlc 2
\begin{ex}%[2H5H2-2]
	Cho đường thẳng $d\colon \dfrac{x-2}{-1}=\dfrac{y-1}{2}=\dfrac{z+3}{1}$. Vectơ nào dưới đây là vectơ chỉ phương của đường thẳng $d$
	\choice
	{$\overrightarrow{u}_1=(2 ; 1 ;-3)$}
	{$\overrightarrow{u}_2=(-2 ;-1 ; 3)$}
	{\True $\overrightarrow{u}_3=(-1 ; 2 ; 1)$}
	{$\overrightarrow{u}_4=(-1 ; 2 ;-1)$}
	\loigiai{
		$d$ có vectơ chỉ phương là $\overrightarrow{u}_3=(-1 ; 2 ; 1)$.
	}
\end{ex}
%nlc 3
\begin{ex}%[2H5H1-6]
	Cho hai mặt phẳng $(P)\colon 2 x-y-z-3=0$ và $(Q)\colon x-z-2=0$. Góc giữa hai mặt phẳng $(P)$ và  $(Q)$ bằng
	\choice
	{$30^{\circ}$}
	{$45^{\circ}$}
	{\True $60^{\circ}$}
	{$90^{\circ}$}
	\loigiai{
		$(P)$ có vectơ pháp tuyến là $\overrightarrow{n}_1=(2 ; -1 ; -1)$.\\
		$(Q)$ có vectơ pháp tuyến là $\overrightarrow{n}_2=(1 ; 0 ; -1)$.\\
		Gọi $\varphi$ là góc giữa hai mặt phẳng $(P)$ và  $(Q)$ ta có $$\cos\varphi=\dfrac{\left| 2\cdot 1-1\cdot 0+ 1\cdot 1\right| }{\sqrt{2^2+(-1)^2+(-1)^2}\cdot\sqrt{1^2+0^2+(-1)^2}}=\dfrac{\sqrt{3}}{2}\Rightarrow \varphi=60^{\circ}.$$
	}
\end{ex}
%nlc 4
\begin{ex}%[2H5H1-6]
	Góc giữa đường thẳng $d\colon \dfrac{x}{1}=\dfrac{y-1}{2}=\dfrac{z}{1}$ và mặt phẳng $(\alpha)\colon 4x+3y+5z-4=0$ bằng
	\choice
	{$30^{\circ}$}
	{$120^{\circ}$}
	{\True $60^{\circ}$}
	{$150^{\circ}$}
	\loigiai{
		$d$ có véc-tơ chỉ phương là $\vec{u}=(1;2;1)$.\\
		$(\alpha)$ có véc-tơ pháp tuyến là $\vec{n}=(4;3;5)$.\\
		$\sin(d,(\alpha))=\dfrac{|\vec{n}\cdot \vec{u}|}{|\vec{n}||\vec{u}|}=\dfrac{|1\cdot 4+2\cdot 3+1\cdot 5|}{\sqrt{1^2+2^2+1^1}\sqrt{4^2+3^2+5^2}}=\dfrac{\sqrt{3}}{2} \Rightarrow (d,(\alpha))=60^{\circ}$.
	}
\end{ex}
%nlc 5
\begin{ex}%[2H5N3-2]
	Trong không gian với hệ toa độ $Oxyz$, phương trình nào sau đây là phương trình mặt cầu?
	\choice
	{$x^{2}+y^{2}+z^{2}+4 x-2 y+6 z+15=0$}
	{$x^{2}+y^{2}+2 z^{2}+4 x-2 y+6 z+5=0$}
	{\True $x^{2}+y^{2}+z^{2}+4 x-2 y+6 z+15=0$}
	{$x^{2}+y^{2}+z^{2}+4 x-2 x y+6 z+5=0$}
	\loigiai{
		\begin{itemize}
			\item $ x^{2}+y^{2}+z^{2}+4 x-2 y+6 z+15=0 \Leftrightarrow (x+2)^2+(y-1)^2+(z+3)^2=-1. $ Đây không phải là phương trình mặt cầu.
			\item $x^{2}+y^{2}+2 z^{2}+4 x-2 y+6 z+5=0$. Đây không phải là phương trình mặt cầu do hệ số $ z^2 $ bằng $ 2 $ khác hệ số của $ x^2 $.
			\item $x^{2}+y^{2}+z^{2}+4 x-2 y+6 z+15=0 \Leftrightarrow (x+2)^2+(y-1)^2+(z+3)^2=9$. Đây là phương trình mặt cầu.
			\item  $x^{2}+y^{2}+z^{2}+4 x-2 x y+6 z+5=0$. Đây không phải là phương trình mặt cầu do hệ số $ xy $ bằng $ -2 $.
		\end{itemize}
	}
\end{ex}
%nlc 6
\begin{ex}%[2H5N3-2]
	Cho mặt cầu $(S)\colon (x+1)^2+(y-2)^2+(z-1)^2=9$. 
	Toạ tâm $I$ và bán kính $R$ của $(S)$ là
	\choice
	{\True $I(-1 ; 2 ; 1)$ và $R=3$}
	{$I(1 ;-2 ;-1)$ và $R=3$}
	{$I(-1 ; 2 ; 1)$ và $R=9$}
	{$I(1 ;-2 ;-1)$ và $R=9$}
	\loigiai{
		Mặt cầu $(S)\colon (x+1)^2+(y-2)^2+(z-1)^2=9$
		có tâm $I(-1 ; 2 ; 1)$ và bán kính $R=3$.
	}
\end{ex}
%nlc 7
\begin{ex}%[2H5H3-3]
	Trong không gian với hệ tọa độ $Oxyz$. Cho mặt cầu
	$(S)\colon x^{2}+y^{2}+z^{2}-2x+4y-4z-m=0$ có bán kính $R=5$. Giá trị $m$ là
	\choice
	{$m=-4$}
	{$m=-16$}
	{\True $m=16$}
	{$m=4$}
	\loigiai{
		Ta có $ x^{2}+y^{2}+z^{2}-2 x+4 y-4 z-m=0 \Leftrightarrow (x-1)^2+(y+2)^2+(z-2)^2=m+9. $
		Bán kính mặt cầu bằng $ 5 $ nên $ m+9=25 \Leftrightarrow m=16 $.
	}
\end{ex}
%nlc 8
\begin{ex}%[2H5H3-3]
	Trong không gian $ Oxyz $, mặt cầu  tâm $I(-1 ; 2 ; 3)$ và tiếp xúc với mặt phẳng $(P)\colon 2x-y-2z+1=0$ có phương trình
	\choice
	{\True $(x+1)^{2}+(y-2)^{2}+(z-3)^{2}=9$}
	{$(x+1)^{2}+(y-2)^{2}+(z-3)^{2}=3$}
	{$(x+1)^{2}+(y-2)^{2}+(z-3)^{2}=2$}
	{$(x+1)^{2}+(y-2)^{2}+(z-3)^{2}=4$}
	\loigiai{
		Bán kính mặt cầu $ R=\mathrm{d}(I,(P))=\dfrac{|2 \cdot (-1)-2-2 \cdot 3+1|}{\sqrt{2^2+(-1)^2+2^2}} =3.$\\
		Phương trình mặt cầu tâm $ I $ tiếp xúc với $ (P) $ là
		$ (x+1)^2+(y-2)^2+(z-3)^2=9. $
	}
	
\end{ex}
%nlc 9
\begin{ex}%[2D5N1-1]
	Cho $A$, $B$ là các biến cố của một phép thử $T$. Biết rằng $\mathrm{P}(B)>0$, xác suất của biến cố $A$ với điều kiện biến cố $B$ đã xảy ra được tính theo công thức nào sau đây?
	\choice
	{$\mathrm{P}(A|B)=\dfrac{\mathrm{P}(A)}{\mathrm{P}(B)}$}
	{$\mathrm{P}(A|B)=\dfrac{\mathrm{P}(A)}{\mathrm{P}(A B)}$}
	{\True $\mathrm{P}(A| B)=\dfrac{\mathrm{P}(A B)}{\mathrm{P}(B)}$}
	{$\mathrm{P}(A|B)=\dfrac{\mathrm{P}(A B)}{\mathrm{P}(A) \cdot \mathrm{P}(B)}$}
	\loigiai{ Dựa theo công thức tính xác suất biến cố $A$ với điều kiện $B$ thì $\mathrm{P}(A| B)=\dfrac{\mathrm{P}(A B)}{\mathrm{P}(B)}$ là đáp án đúng.
		
	}
\end{ex}
%nlc 10
\begin{ex}%[2D5H1-2]
	\immini[thm]{Cho sơ đồ hình cây như hình bên. Xác suất điều kiện $\mathrm{P}(A \mid B)$ là
		\choice
		{\True$\dfrac{7}{31}$}
		{$0\text{,}7$}
		{$\dfrac{7}{50}$}
		{$0\text{,}48$}}
	{\begin{tikzpicture}[scale=.2,>=stealth]
			\tikzstyle{block} = [rectangle, draw, fill=blue!10\text{,} rounded corners, text centered, text width = 10em, minimum height = 2em]
			\node (c1) {};
			\node (c2)[above right = 1.5cm of c1] {$A$};
			\node at (0.5,5){\fbox{$0\text{,}2$}};
			\node at (0.5,-5){\fbox{$0\text{,}8$}};
			\node (c3) [below right= 1.5cm of c1]{$\overline{A}$};
			\node at (12,11.5){\fbox{$0\text{,}7$}};
			\node (c4) at (21.5, 12){$B$};
			\node (c5) at (21.5, 2){$\overline{B}$};
			\node at (12,3){\fbox{$0\text{,}3$}};
			\node (c6) at (21.5, -4){$B$};
			\node at (12,-4){\fbox{$0\text{,}6$}};
			\node (c7) at (21.5, -14){$\overline{B}$};
			\node at (12,-13){\fbox{$0\text{,}4$}};
			\draw[->] (c1.east) -- (c2.west);
			\draw[->] (c1.east) -- (c3.west);
			\draw[->] (c2.east) -- (c4.west);
			\draw[->] (c2.east) -- (c5.west);
			\draw[->] (c3.east) -- (c6.west);
			\draw[->] (c3.east) -- (c7.west);
	\end{tikzpicture}}
	\loigiai{
		\begin{itemize}
			\item Ta có $\mathrm{P}(AB)=0\text{,}2\cdot 0\text{,}7=0\text{,}14$.
			\item Mặt khác $\mathrm{P}(B)=0\text{,}2\cdot 0\text{,}7+0\text{,}8\cdot 0\text{,}6=0\text{,}62$.
			\item $\mathrm{P}(A\mid B)=\dfrac{\mathrm{P}(AB)}{\mathrm{P}(B)}=\dfrac{0\text{,}14}{0\text{,}62}=\dfrac{7}{31}$.
		\end{itemize}	
	}
\end{ex}
%nlc 11
\begin{ex}%[2D5V2-3]
	Cho $A$, $B$ là các biến cố thỏa mãn $\mathrm{P}(\overline{A} \,\overline{B})=0{,}35$, $\mathrm{P}(A)=0{,}25$, $\mathrm{P}(B)=0{,}6$. Giá trị của $\mathrm{P}(A|B)$ bằng
	\choice
	{$\dfrac{1}{5}$}
	{\True$\dfrac{1}{3}$}
	{$\dfrac{7}{15}$}
	{$\dfrac{2}{3}$}
	\loigiai{Ta có $\mathrm{P}\left( \overline{A}\,\overline{B}  \right) = \mathrm{P}\left( {\overline A } \right)\mathrm{P}\left( {\overline B |\overline A } \right) \Rightarrow \mathrm{P}\left( {\overline B |\overline A } \right) = \dfrac{{\mathrm{P}\left( \overline{A}\,\overline{B} \right)}}{{\mathrm{P}\left( {\overline A } \right)}} = \dfrac{{0{,}35}}{{0{,}75}} = \dfrac{7}{{15}}.$ \\
		Suy ra $\mathrm{P}\left( {B|\overline A } \right) = 1 - \dfrac{7}{{15}} = \dfrac{8}{{15}}$.\\
		Theo công thức xác suất toàn phần, ta có
		$$\begin{array}{l}
			\mathrm{P}\left( B \right) = \mathrm{P}\left( A \right)\mathrm{P}\left( {B|A} \right) +\mathrm{P}\left( {\overline A } \right) \mathrm{P}\left( {B|\overline A } \right)\\
			\Rightarrow \mathrm{P}\left( {B|A} \right) = \dfrac{{\mathrm{P}\left( B \right) - \mathrm{P}\left( {B|\overline A } \right)\mathrm{P}\left( {\overline A } \right)}}{{\mathrm{P}\left( A \right)}} = \dfrac{{0{,}6 - \dfrac{8}{{15}} \cdot 0{,}75}}{{0{,}25}} = 0{,}8.
		\end{array}.$$
		Theo công thức Bayes, ta được
		$$\mathrm{P}\left( {A|B} \right) = \dfrac{{\mathrm{P}\left( A \right)\mathrm{P}\left( {B|A} \right)}}{{\mathrm{P}\left( B \right)}} = \dfrac{{0{,}25 \cdot 0{,}8}}{{0{,}6}} = \dfrac{1}{3}.$$}
\end{ex}
%nlc 12
\begin{ex}%[2D5V2-3]
	Một bệnh viện có hai phòng khám là phòng A và phòng B với khả năng lựa chọn của bệnh nhân là như nhau. Tỉ lệ bệnh nhân nam có ở phòng A và phòng B lần lượt là $60\%$ và $40\%$. Một người bệnh được chọn ngẫu nhiêu từ hai phòng khám và biết người này là nam, xác suất để người bệnh được chọn đến từ phòng A là
	\choice
	{\True $0{,}6$}
	{$0{,}5$}
	{$0{,}4$}
	{$0{,}3$}
	\loigiai{Một người bệnh được chọn ngẫu nhiên từ hai phòng khám.\\
		Gọi $X$ là biến cố \lq \lq Người đó đến từ phòng khám A\rq \rq \, và $Y$, $\overline{Y}$ lần lượt là biến cố \lq \lq Người đó là nam\rq \rq \; và \lq \lq Người đó không là nam\rq \rq.\\
		Ta có sơ đồ hình cây sau\\
		\begin{tikzpicture}[>=stealth]
			%Khung 1
			\draw (-2,-1) rectangle (2.2,0);
			\draw (0.1,-0.5) node{Bệnh nhân được chọn} ;
			%Mui ten 1,2
			\draw [->] (2.2,-0.5)--(3.8,1.6) node[pos=0.5,sloped,above]{$0{,}5$};
			\draw [->] (2.2,-0.5)--(3.8,-2.6) node[pos=0.5,sloped,below]{$0{,}5$};
			%Khung 2.1
			\draw (3.8,1.1) rectangle (5.1,2.1);
			\draw (8.9/2,1.6) node{$X$} ;
			%Khung 2.2
			\draw (3.8,-2.1) rectangle (5.1,-3.1);
			\draw (8.9/2,-2.6) node{$\overline{X}$} ;
			%Mui ten 3,4
			\draw [->] (5.1,1.6)--(6.5,2.6) node[pos=0.5,sloped,above]{$0{,}6$};
			\draw [->] (5.1,1.6)--(6.5,0.6) node[pos=0.5,sloped,below]{$0{,}4$};
			%Mui ten 5,6
			\draw [->] (5.1,-2.6)--(6.5,-1.6) node[pos=0.5,sloped,above]{$0{,}4$};
			\draw [->] (5.1,-2.6)--(6.5,-3.6) node[pos=0.5,sloped,below]{$0{,}6$};
			%Khung 3.1
			\draw (6.5,2.2) rectangle (7.7,3.2);
			\draw (7.1,5.4/2) node{$Y$} ;
			%Khung 3.2
			\draw (6.5,1.2) rectangle (7.7,0.2);
			\draw (7.1,1.4/2) node{$\overline{Y}$} ;
			%Khung 3.3
			\draw (6.5,-1.1) rectangle (7.7,-2.1);
			\draw (7.1,-3.2/2) node{$Y$} ;
			%Khung 3.3
			\draw (6.5,-2.9) rectangle (7.7,-3.9);
			\draw (7.1,-3.4) node{$\overline{Y}$} ;
			%Kết quả
			\draw (9.5,3.7) node{\textbf{Kết quả}};	
			\draw (9.5,2.7) node{$XY$};
			\draw (9.5,0.7) node{$X \overline{Y}$};
			\draw (9.5,-1.6) node{$\overline{X}Y$};
			\draw (9.5,-3.4) node{$\overline{X}\overline{Y}$};
			%Xác suất
			\draw (12.5,3.7) node{\textbf{Xác suất}};	
			\draw (12.5,2.7) node{$0{,}3$};
			\draw (12.5,0.7) node{$0{,}2$};
			\draw (12.5,-1.6) node{$0{,}2$};
			\draw (12.5,-3.4) node{$0{,}3$};	
		\end{tikzpicture}
		
		Theo công thức Bayes, ta có $$\mathrm{P}(X|Y)=\dfrac{\mathrm{P}(X)\mathrm{P}(Y|X)}{\mathrm{P}(X)\mathrm{P}(Y|X)+\mathrm{P}(\overline{X})\mathrm{P}(Y|\overline{X})}=\dfrac{0{,}3}{0{,}3+0{,}2}=0{,}6.$$
		Vậy với một người bệnh được chọn ngẫu nhiêu từ hai phòng khám và biết người này là nam, xác suất để người đó đến từ phòng A là $0{,}6$.}
\end{ex}
\Closesolutionfile{ans}
\indapan{6}{ans/ans-2-OHK2-De3}

\cauds
\Opensolutionfile{ans}[ans/ans-2-OHK2-De3-DS]
%ds 1
\begin{ex}%[2H5N3-2]
	Trong không gian $Oxyz$, phương trình mặt cầu $$(S) \colon (x-1)^{2} + (y+2)^{2} + (z-1)^{2} =~25.$$
	Các mệnh đề sau đúng hay sai?
	\choiceTF
	{\True Tâm của mặt cầu $(S)$ là $I(1,-2,1)$}
	{\True Bán kính của mặt cầu $(S)$ là $R = 5$}
	{Điểm $M(5,1,-1)$ nằm trên mặt cầu $(S)$}
	{Đường kính của mặt cầu $(S)$ là $25$}
	\loigiai{
		\begin{itemchoice}
			\itemch Đúng.
			\itemch Đúng. $R^{2} = 25 \Rightarrow R = 5$.
			\itemch Sai. Thay điểm $M(5,1,-1)$ vào phương trình mặt cầu\\
			$\Rightarrow (5-1)^{2} + (1+2)^{2} + (-1-1)^{2} = 29 \ne 25$.\\
			Suy ra, điểm $M \not\in (S)$.
			\itemch Đường kính mặt cầu là $2R = 10$.
		\end{itemchoice}
	}
\end{ex}
%ds 2
\begin{ex}
	Tại một tiệm nhỏ bán điểm tâm sáng của chị Loan, chị Loan chỉ bán bún riêu hoặc hủ tiếu. Nếu hôm nay chị bán bún riêu thì xác suất để hôm sau chị bán hủ tiếu là $0{,}2$. Nếu hôm nay chị bán hủ tiếu thì xác suất để hôm sau chị bán bún riêu là $0{,}6$. Xét một tuần mà thứ Ba, chị Loan bán bún riêu. Tính xác suất để thứ Năm trong tuần đó, chị Loan bán hủ tiếu.
	\choiceTF
	{$0{,}2 \cdot 0{,}4 + 0{,6} \cdot 0{,}4$}
	{\True $0{,}2 \cdot 0{,}4 + 0{,8} \cdot 0{,}2$}
	{$0{,}6 \cdot 0{,}4 + 0{,6} \cdot 0{,}2$}
	{\True $0{,8} \cdot 0{,}2 + 0{,}2 \cdot 0{,}4$}
	\loigiai{
		\begin{itemchoice}
			\itemch Sai.
			\itemch Đúng.\\
			+ Gọi $A$ là biến cố: \lq\lq Thứ Tư, chị Loan bán hủ tiếu\rq\rq. Khi đó, $\mathrm{P}(A) = 0{,}2$ và $\mathrm{P}(\overline{A}) = 1 - 0{,}2 = 0{,}8$.\\
			+ Gọi $B$ là biến cố:  \lq\lq Thứ Năm, chị Loan bán hủ tiếu\rq\rq.\\
			+ Xác suất để chị Loan thứ Tư bán hủ tiếu và thứ Năm cũng bán hủ tiếu là $\mathrm{P}(B|A) = 1 - 0{,}6 = 0{,}4$.\\
			+ Xác suất để chị Loan thứ Tư bán bún riêu và thứ Năm cũng bán hủ tiếu là $\mathrm{P}(B|\overline{A}) = 0{,}2$.\\
			+ Xác suất để ngày thứ Năm chị Loan bán hủ tiếu là\\
			$\Rightarrow \mathrm{P}(B) = \mathrm{P}(A) \cdot \mathrm{P}(B|A) + \mathrm{P}(\overline{A}) \cdot \mathrm{P}(B|\overline{A}) = 0{,}2 \cdot 0{,}4 + 0{,}8 \cdot 0{,}2$.
			\itemch Sai.
			\itemch Đúng.\\
			+ Gọi $A$ là biến cố: \lq\lq Thứ Tư, chị Loan bán bún riêu\rq\rq. Khi đó, $\mathrm{P}(A) = 1 - 0{,}2 = 0{,}8$ và $\mathrm{P}(\overline{A}) = 0{,}2$.\\
			+ Gọi $B$ là biến cố: \lq\lq Thứ Năm, chị Loan bán hủ tiếu\rq\rq.\\
			+ Xác suất để chị Loan thứ Tư bán hủ tiếu và thứ Năm cũng bán hủ tiếu là $\mathrm{P}(B|A) = 0{,}2$.\\
			+ Xác suất để chị Loan thứ Tư bán bún riêu và thứ Năm cũng bán hủ tiếu là $\mathrm{P}(B|\overline{A}) = 1 - 0{,}6 = 0{,}4$.\\
			+ Xác suất để ngày thứ Năm chị Loan bán hủ tiếu là\\
			$\Rightarrow \mathrm{P}(B) = \mathrm{P}(A) \cdot \mathrm{P}(B|A) + \mathrm{P}(\overline{A}) \cdot \mathrm{P}(B|\overline{A}) = 0{,}8 \cdot 0{,}2 + 0{,}2 \cdot 0{,}4$.
		\end{itemchoice}
	}
\end{ex}
%ds 3
\begin{ex}%[2H5H2-7]
	Trong không gian $Oxyz$, cho đường thẳng $\Delta$ và mặt phẳng $(\alpha)$ có phương trình
	$$\Delta \colon \heva{&x = 1 + t\\&y = 2 - t\\&z = -1 + 2t}\qquad (\alpha) \colon x + 2y - 3z + 5 = 0$$
	Các mệnh đề sau đúng hay sai?
	\choiceTF
	{\True Đường thẳng $\Delta$ và mặt phẳng $(\alpha)$ cắt nhau}
	{Đường thẳng $\Delta$ và mặt phẳng $(\alpha)$ vuông góc với nhau}
	{Đường thẳng $\Delta$ và mặt phẳng $(\alpha)$ song song với nhau}
	{\True Góc giữa đường thẳng $\Delta$ và mặt phẳng $(\alpha)$ xấp xỉ $49{,}5^{\circ}$}
	\loigiai{
		\begin{itemchoice}
			\itemch Đúng. Vì $\vec{u}_{\Delta} \cdot \vec{n}_{(\alpha)} = -7 \ne 0$, tức là 2 vec-tơ này không vuông góc. Suy ra đường thẳng và mặt phẳng không song song.
			\itemch Sai. Vì $\dfrac{1}{1} \ne \dfrac{-1}{2} \ne \dfrac{2}{-3}$. Hai vec-tơ không cùng phương. Đường thẳng và mặt phẳng không vuông góc nhau.
			\itemch Sai. Vì $\vec{u}_{\Delta} \cdot \vec{n}_{(\alpha)} = -7 \ne 0$, nên chúng không song song nhau.
			\itemch Đúng. Vì\\
			$\sin(\Delta, (\alpha)) = \dfrac{|1 \cdot 1 + (-2)\cdot 2 + 2 \cdot (-3)|}{\sqrt{1^{2} + (-1)^{2} + 2^{2}} \cdot \sqrt{1^{2} + 2^{2} + (-3)^{2}}} \approx 0{,}76$.\\
			Suy ra góc $(\Delta, (\alpha)) \approx 49{,}5^{\circ}$.
		\end{itemchoice}
	}
\end{ex}
%ds 4
\begin{ex}%[2D5V1-2]
	Một nghiên cứu về an toàn giao thông muốn tìm hiểu về mối quan hệ giữa việc thắt dây an toàn khi lái xe và nguy cơ tử vong của người lái xe khi xảy ra tai nạn giao thông. Giả sử đã xem xét $845814$ vụ tai nạn giao thông ô tô và việc thắt dây an toàn của người lái xe khi xảy ra tai nạn giao thông. Kết quả cho thấy:
	\begin{itemize}
		\item Trong số những người lái xe có thắt dây an toàn, có $515$ người tử vong và $621188$ người sống sót;
		\item Trong số những người lái xe không thắt dây an toàn, có $1605$ người tử vong và $222506$ người sống sót.
	\end{itemize}
	Các mệnh đề sau là đúng hay sai?
	\choiceTF
	{Không gian mẫu $n(\Omega) = 843694$}
	{Xác suất người lái xe đó tử vong khi xảy ra tai nạn trong trường hợp không thắt dây an toàn là $\dfrac{1605}{222506}$}
	{\True Xác suất người lái xe đó tử vong khi xảy ra tai nạn trong trường hợp có thắt dây an toàn là $\dfrac{515}{621703}$}
	{\True Tỉ số giữa xác suất của người lái xe tử vong khi không thắt dây ăn toàn và người tử vong khi có thắt dây an toàn xấp xỉ $8{,}6$ lần}
	\loigiai{
		+ Không gian mẫu $n(\Omega) = 845814$.\\
		+ Gọi $A$ là biến cố: \lq\lq Người lái xe tử vong khi xảy ra tai nạn giao thông\rq\rq;\\
		+ Gọi $B$ là biến cố: \lq\lq Người lái xe không thắt dây an toàn\rq\rq.\\
		Khi đó $AB$ là biến cố: \lq\lq Người lái xe tử vong không thắt dây an toàn khi xảy ra tai nạn giao thông\rq\rq.\\
		+ Xác suất tử vong, không thắt dây an toàn khi xảy ra tai nạn giao thông:\\
		$\mathrm{P}(A|B) = \dfrac{\mathrm{P}(AB)}{\mathrm{P}(B)} = \dfrac{n(AB)}{n(B)} = \dfrac{1605}{1605 + 222506} = \dfrac{1605}{224111}$.\\
		+ Xác suất tử vong, có thắt dây an toàn khi xảy ra tai nạn giao thông:\\
		$\mathrm{P}(A|\overline{B}) = \dfrac{\mathrm{P}(A\overline{B})}{\mathrm{P}(\overline{B})} = \dfrac{n(A\overline{B})}{n(\overline{B})} = \dfrac{515}{515 + 621188} = \dfrac{515}{621703}$.\\
		+ Tỉ số giữa xác suất của người lái xe tử vong khi không thắt dây ăn toàn và người tử vong khi có thắt dây an toàn
		$\dfrac{\mathrm{P}(A|B)}{\mathrm{P}(A|\overline{B})} \approx 8{,}6$.
		\begin{itemchoice}
			\itemch Sai. Không gian mẫu $n(\Omega) = 545814$.
			\itemch Sai. $\mathrm{P}(A|B) = \dfrac{1605}{224111}$.
			\itemch Đúng. $\mathrm{P}(A|\overline{B}) = \dfrac{151}{621703}$.
			\itemch Đúng. $\dfrac{\mathrm{P}(A|B)}{\mathrm{P}(A|\overline{B})} \approx 8{,}6$.
		\end{itemchoice}
	}
\end{ex}
\Closesolutionfile{ans}
\indapan{3}{ans/ans-2-OHK2-De3-DS}

\Opensolutionfile{ans}[ans/ans-2-OHK2-De3-KQ]
\caukq
%tln 1
\begin{ex}%[2H5H1-3]
	Trong không gian với hệ trục toạ độ $Oxyz$, cho ba điểm $A(2;0;0)$; $B(0;-1;0)$; $C(0;0;3)$. Phương trình mặt phẳng $(ABC)$ có dạng $ax+by+cz-6=0$. Khi đó $a+b+c$ bằng bao nhiêu?
	\shortans{$-1$}
	\loigiai{
		Mặt phẳng $(ABC)$  có phương trình đoạn chắn là 
		$$\dfrac{x}{2}+ \dfrac{y}{-1} + \dfrac{z}{3}=1 \Leftrightarrow 3x - 6y + 2z -6=0.$$
		Vậy $a+b+c=3-6+2=-1$.
	}
\end{ex}
%tln 2
\begin{ex}%[2H5H1-6]
	Cho hai mặt phẳng $(P)\colon x-y-6=0$ và $(Q)$. Biết rằng điểm $H(2 ;-1 ;-2)$ là hình chiếu vuông góc của gốc toạ độ $O(0;0;0)$ xuống mặt phẳng $(Q)$. Góc giữa mặt phẳng $(P)$ và mặt phẳng $(Q)$ là $a^\circ$. Khi đó $a$ bằng bao nhiêu?
	\shortans{$45$}
	\loigiai{		
		$(P)$ có vectơ pháp tuyến là $\overrightarrow{n}_1=(1 ; -1 ; 0)$, 	$(Q)$ có vectơ pháp tuyến là $ \overrightarrow {OH}=(2 ; -1 ; -2)$.\\ 
		Gọi $\varphi$ là góc giữa hai mặt phẳng $(P)$ và  $(Q)$ ta có $$\cos\varphi=\dfrac{|2\cdot 1+1\cdot 1-2\cdot 0|}{\sqrt{1^2+(-1)^2+0^2}\cdot\sqrt{2^2+(-1)^2+(-2)^2}}=\dfrac{\sqrt{2}}{2}\Rightarrow \varphi=45^{\circ}.$$
	}
\end{ex}
%tln 3
\begin{ex}%[2H5N3-2]
	Trong không gian với hệ tọa độ $ Oxyz $ cho mặt cầu 
	$(S)\colon x^{2}+y^{2}+z^{2}+4 x-8 y+2 z-4=0$. Bán kính $R$ của $(S)$ là
	\shortans{$5$}
	\loigiai{
		Ta có $ x^{2}+y^{2}+z^{2}+4 x-8 y+2 z-4=0 \Leftrightarrow (x+2)^2+(y-4)^2+(z+1)^2=25. $\\
		$ (S) $ có tâm $ I(-2;4;-1) $ có bán kính $R=5$.
	}
\end{ex}
%tln 4
\begin{ex}%[2H5C3-4]
	Trong không gian với hệ tọa độ $Oxyz$, đài kiểm soát không lưu sân bay có tọa độ $O(0;0;0)$, mỗi đơn vị trên trục ứng với $1$ km. Máy bay bay trong phạm vi cách đài kiểm soát $417$ km sẽ hiển thị trên màn hình ra đa. Một máy bay đang ở vị trí $A\left(-688;-185;8\right)$, chuyển động theo đường thẳng $d$ có véc-tơ chỉ phương là $\vec{u}=(91;75;0)$ và hướng về đài kiểm soát không lưu.
	\begin{center}
		\begin{tikzpicture}
			\def\r{3}
			\def \x{3}  %bán kính trục lớn elip
			\def \y{1.2} 
			\coordinate (o1) at (-\r,0);
			\coordinate (o2) at (\r,0);
			\coordinate (O) at (0,0);
			
			% ve mp
			\coordinate (P) at (-5,2);
			\coordinate (Q) at ($(P)+(240:6)$);
			\coordinate (R) at ($(Q)+(12,0)$);
			\coordinate (S) at ($(R)+(60:6)$);
			\fill[cyan!20] (P)--(Q)--(R)--(S)--cycle;
			
			
			\draw (O) circle(\r); 
			\path[fill=orange!20] (O) circle(\r); 
			\fill[orange!80] (o2) arc (0:-180:\x cm and \y cm) arc (180:360:\x cm and \x cm);
			\draw[dashed] (o2) arc (0:180:\x cm and \y cm);
			\draw (o2) arc (0:-180:\x cm and \y cm);
			
			% truc oz
			\coordinate (Oz) at (0,4.5);
			\coordinate (Oz') at (0,2.5);
			\coordinate (Oz'') at (0,-3.5);
			\draw[dashed] (Oz'')--(Oz');
			\draw[->] (Oz')--(Oz);
			% tru  ox
			\def\t{240}
			\coordinate (Ox) at ({\x*cos(\t)},{\y*sin(\t)});
			\coordinate (Ox') at ($(O)!2!(Ox)$);
			\coordinate (Ox'') at ($(Ox)!2.2!(O)$);
			\draw[dashed] (O)--(Ox);
			\draw[->] (Ox)--(Ox');
			\draw[dashed] (O)--(Ox'');
			% truc oy
			\coordinate (O4) at (\r,0);
			\coordinate (O5) at (1.5*\r,0);
			\coordinate (O2) at (-\r,0);
			\coordinate (O1) at (-1.5*\r,0);
			\draw (O1)--(O2);
			\draw[dashed] (O2)--(O4);
			\draw[->] (O4)--(O5);
			% duong d
			\coordinate (A) at (-4.52,1.7);
			\coordinate (A3) at (4.5,0.3);
			\coordinate (A1) at ($(A)!0.25!(A3)$);
			\coordinate (A2) at ($(A)!0.82!(A3)$);
			\draw (A)--(A1);
			\draw[dashed] (A1)--(A2);
			\draw (A2)--(A3);
			\foreach \x in {O}{
				\fill (\x) circle(1pt);
			}
			\path 
			node at (O) [above left]{$O$}
			node at (A) [above] {$A$}
			node at (A1) [above]{B}
			node at (A2) [above right]{$C$}
			node at (A3) [above]{$d$}
			node at (Oz) [left]{$z$}
			node at (O5) [below]{$y$}
			node at (Ox')[below]{$x$}
			;
		\end{tikzpicture}
	\end{center}		
	Máy bay bay gần đài kiểm soát không lưu nhất ở vị trí $K(a;b;c)$. Khi đó $a+b+c$ bằng bao nhiêu?
	\shortans{$48$}		
	\loigiai{
		Gọi $(P)$ là mặt phẳng đi qua $O$ và vuông góc với đường thẳng $d$. \\
		Ta có $(P) \perp d$ nên nhận véc-tơ chỉ phương của $d$ là một véc-tơ pháp tuyến.\\
		Vậy $(P)$ đi qua $O(0;0;0)$ và có véc-tơ pháp tuyến $\vec{n}_{P}=\vec{u}_d=(91;75;0)$ nên có phương trình $91\cdot (x-0)+75\cdot (y-0)+0\cdot (z-0)=0 \Leftrightarrow 91x+75y=0$.\\
		Gọi $K$ là hình chiếu vuông góc của $O$ lên đường thẳng $d$, ta có $K$ là giao của đường thẳng $d$ và mặt phẳng $(P)$. Khi đó tọa độ của $K$ là nghiệm của hệ phương trình 
		\begin{eqnarray*}
			\heva{&x=-688 +91t \\ &y=-185+75t\\ &z=8 \\& 91x+75y=0} 
			\Leftrightarrow
			\heva{&x=-688 +91t \\ &y=-185+75t\\ &z=8 \\& 91 \cdot(-688 +91t)+75 \cdot (-185+75t)=0}  
			\Leftrightarrow \heva{& t=\dfrac{11}{2} \\& x=-\dfrac{375}{2}\\&y=\dfrac{455}{2}\\&z=8.}
		\end{eqnarray*}
		Vậy $K\left(-\dfrac{375}{2};\dfrac{455}{2};8\right)$. Tọa độ của vị trí mà máy bay bay gần đài kiểm soát không lưu nhất chính là $K\left(-\dfrac{375}{2};\dfrac{455}{2};8\right)$.\\						
	}
\end{ex}
%tln 5
\begin{ex}%[2D5V2-3]
	Một bệnh viện đang xét nghiệm cho một số bệnh nhân để xác định liệu họ có nhiễm virus $X$ hay không. Xác suất để một bệnh nhân bị nhiễm virus $X$ là $0{,}05$. Khi xét nghiệm, nếu một bệnh nhân bị nhiễm thì xác suất để kết quả xét nghiệm dương tính là $0{,}95$. Nếu một bệnh nhân không bị nhiễm thì xác suất để kết quả xét nghiệm âm tính là $0{,}98$. Một bệnh nhân được chọn ngẫu nhiên và có kết quả xét nghiệm dương tính. Xác suất để bệnh nhân đó thực sự bị nhiễm virus $X$ bằng bao nhiêu? (làm tròn kết quả đến hàng phần mười) 
	\shortans{$0{,}7$}	
	\loigiai{Một bệnh nhân đến một bệnh viên để xét nghiệm.\\
		Gọi $A$ là biến cố \lq \lq Bệnh nhân bị nhiễm virus $X$\rq \rq \, và $B$, $\overline{B}$ lần lượt là biến cố \lq \lq Kết quả xét nghiệm dương tính\rq \rq \; và \lq \lq Kết quả xét nghiệm âm tính\rq \rq.\\
		Ta xét sơ đồ hình cây như sau\\
		\begin{tikzpicture}[>=stealth]
			%Khung 1
			\draw (-3.5,-1) rectangle (2.2,0);
			\draw (-1.3/2,-0.5) node{Bệnh nhân được xét nghiệm} ;
			%Mui ten 1,2
			\draw [->] (2.2,-0.5)--(3.8,1.6) node[pos=0.5,sloped,above]{$0{,}05$};
			\draw [->] (2.2,-0.5)--(3.8,-2.6) node[pos=0.5,sloped,below]{$0{,}95$};
			%Khung 2.1
			\draw (3.8,1.1) rectangle (5.1,2.1);
			\draw (8.9/2,1.6) node{$A$} ;
			%Khung 2.2
			\draw (3.8,-2.1) rectangle (5.1,-3.1);
			\draw (8.9/2,-2.6) node{$\overline{A}$} ;
			%Mui ten 3,4
			\draw [->] (5.1,1.6)--(6.5,2.6) node[pos=0.5,sloped,above]{$0{,}95$};
			\draw [->] (5.1,1.6)--(6.5,0.6) node[pos=0.5,sloped,below]{$0{,}05$};
			%Mui ten 5,6
			\draw [->] (5.1,-2.6)--(6.5,-1.6) node[pos=0.5,sloped,above]{$0{,}02$};
			\draw [->] (5.1,-2.6)--(6.5,-3.6) node[pos=0.5,sloped,below]{$0{,}98$};
			%Khung 3.1
			\draw (6.5,2.2) rectangle (7.7,3.2);
			\draw (7.1,5.4/2) node{$B$} ;
			%Khung 3.2
			\draw (6.5,1.2) rectangle (7.7,0.2);
			\draw (7.1,1.4/2) node{$\overline{B}$} ;
			%Khung 3.3
			\draw (6.5,-1.1) rectangle (7.7,-2.1);
			\draw (7.1,-3.2/2) node{$B$} ;
			%Khung 3.3
			\draw (6.5,-2.9) rectangle (7.7,-3.9);
			\draw (7.1,-3.4) node{$\overline{B}$} ;
			%Kết quả
			\draw (9.5,3.7) node{\textbf{Kết quả}};	
			\draw (9.5,2.7) node{$AB$};
			\draw (9.5,0.7) node{$A\overline{B}$};
			\draw (9.5,-1.6) node{$\overline{A}B$};
			\draw (9.5,-3.4) node{$\overline{A}\overline{B}$};
			%Xác suất
			\draw (12.5,3.7) node{\textbf{Xác suất}};	
			\draw (12.5,2.7) node{$0{,}0475$};
			\draw (12.5,0.7) node{$0{,}0025$};
			\draw (12.5,-1.6) node{$0{,}019$};
			\draw (12.5,-3.4) node{$0{,}931$};
		\end{tikzpicture}
		
		Theo công thức Bayes, ta có $$\mathrm{P}(A|B)=\dfrac{\mathrm{P}(A)\mathrm{P}(B|A)}{\mathrm{P}(A)\mathrm{P}(B|A)+\mathrm{P}(\overline{A})\mathrm{P}(B|\overline{A})}=\dfrac{0{,}0475}{0{,}0475+0{,}019}=\dfrac{5}{7}.$$		 
		Vậy với một bệnh nhân có kết quả xét nghiệm dương tính, xác suất để bệnh nhân đó thực sự bị nhiễm virus $X$ là $\dfrac{5}{7}\approx 0{,}7$.}
\end{ex}
%tln 6
\begin{ex}%[2D5C2-3]
	Một doanh nghiệp có $45 \%$ nhân viên là nữ. Tỉ lệ nhân viên nữ và tỉ lệ nhân viên nam mua bảo hiểm nhân thọ lần lượt là $7 \%$ và $5 \%$. Gặp ngẫu nhiên một nhân viên của doanh nghiệp.
	Biết rằng nhân viên đó có mua bảo hiểm nhân thọ. Xác suất nhân viên đó là nam băng bao nhiêu? (làm tròn kết quả đến hàng phần mười) 
	\shortans{$0{,}5$}	
	\loigiai{
		\begin{itemize}
			\item $A$: \lq\lq Nhân viên có mua bảo hiểm nhân thọ\rq\rq.
			\item $B$: \lq\lq Nhân viên là nữ \rq\rq.
		\end{itemize}
		Do doanh nghiệp có $45\%$ nhân viên là nữ cho nên \[\mathrm{P}(B)=0{,}45 ~\text{và}~ \mathrm{P}(\overline{B})=0{,}55.\]
		Mặt khác tỉ lệ nhân viên nữ và tỉ lệ nhân viên nam mua bảo hiểm nhân thọ lần lượt là $7 \%$ và $5 \%$ cho nên ta có \[\mathrm{P}(A\mid B)=0{,}07~ \text{và}~ \mathrm{P}(A\mid \overline{B})=0{,}05.\]
		Xác suất nhân viên có mua bảo hiểm nhân thọ là 
		\[ \mathrm{P}(A)=\mathrm{P}(B)\mathrm{P}(A\mid B)+\mathrm{P}(\overline{B})\mathrm{P}(A\mid \overline{B})=0{,}45\cdot 0{,}07
		+ 0{,}55\cdot 0{,}05 = 0{,}059.\]		
		Xác suất nhân viên nam mua bảo hiểm nhân thọ
		\[ \mathrm{P}(\overline{B} \mid A)=\dfrac{\mathrm{P}(\overline{B})\mathrm{P}(A\mid \overline{B})}{\mathrm{P}(A)}=\dfrac{0{,}05 \cdot 0{,}55}{0{,}059} =\dfrac{55}{118}\approx 0{,}5.\]
	}
\end{ex}
\Closesolutionfile{ans}
\indapan{6}{ans/ans-2-OHK2-De3-KQ}